\setcounter{table}{3}
\begin{table}[H]
\centering
\caption{Representative perturbation, counterfactual, and corruption-based evaluation studies in the dataset corpus.}
\label{tab:perturbation_methods}
\renewcommand{\arraystretch}{1.2}
\begin{tabularx}{\linewidth}{>{\raggedright\arraybackslash}p{2.5cm} >{\raggedright\arraybackslash}X >{\raggedright\arraybackslash}X}
\hline
\textbf{Reference} & \textbf{Method} & \textbf{Application} \\
\hline
\cite{Zhang2021} 2021         & Input perturbation and attribution comparison   & ECG diagnosis \\
\cite{MorenoSanchez2023} 2023 & Explainable feature ablation                    & Heart-failure survival \\
\cite{singla2023explaining} 2023 & Counterfactual explanation generation        & Medical image analysis \\
\cite{lenis2020domain} 2020   & Counterfactual impact analysis                  & Medical image classifier interpretation \\
\cite{palatnik2021xai} 2021   & Bias-sensitive saliency analysis                & COVID CT classifier auditing \\
\cite{Pertuz2023} 2023        & Saliency-mask comparison                        & Breast lesion detection \\
\cite{Repetto2025EvaluatingTR} 2025 & Corruption-based robustness testing       & Medical image recognition \\
\cite{Patricio2024} 2024      & AOPC/MoRF perturbation evaluation of attribution maps & Medical image classification \\
\cite{sun2023improving} 2023  & Patch perturbation-based evaluation pipeline    & COVID-19 X-ray analysis \\
\cite{jacob2025indian} 2025   & Cross-corpus perturbation-based XAI evaluation  & Wound healing classification \\
\cite{akgundougdu2025explainable} 2025 & Multi-method perturbation (LIME, Grad-CAM, SHAP) & Brain tumor detection \\
\cite{alpsalaz2026explainability} 2026 & Reliability-weighted perturbation evaluation & Cervical cancer classification \\
\hline
\end{tabularx}
\end{table}
